\chapter{Sandbox (Isolated Execution)}

\section{Overview}

OpenClaw can \textbf{run tools inside Docker containers} to reduce the blast radius. Sandbox isolation is \textbf{optional} and controlled by configuration. The Gateway itself remains on the host; only tool execution is isolated in containers.

Config path: \texttt{agents.defaults.sandbox.*}

\section{What Gets Sandboxed}

\begin{center}
\begin{tabular}{ll}
\toprule
\textbf{Sandboxed} & \textbf{Not Sandboxed} \\
\midrule
Tool execution (exec, read, write, edit, etc.) & Gateway process itself \\
Sandbox browser (optional) & Elevated tools (\texttt{tools.elevated}) \\
 & Tools explicitly allowed to run on host \\
\bottomrule
\end{tabular}
\end{center}

\section{Sandbox Modes}

\texttt{agents.defaults.sandbox.mode}:

\begin{center}
\begin{tabular}{ll}
\toprule
\textbf{Mode} & \textbf{Behavior} \\
\midrule
\texttt{off} & No sandbox \\
\texttt{non-main} & Sandbox for non-main sessions only (default) \\
\texttt{all} & All sessions run in sandbox \\
\bottomrule
\end{tabular}
\end{center}

\section{Scope}

\texttt{agents.defaults.sandbox.scope}:

\begin{itemize}
  \item \texttt{session} (default): One container per session
  \item \texttt{agent}: One container per agent
  \item \texttt{shared}: All sandboxed sessions share one container
\end{itemize}

\section{Workspace Access}

\texttt{agents.defaults.sandbox.workspaceAccess}:

\begin{center}
\begin{tabular}{ll}
\toprule
\textbf{Value} & \textbf{Behavior} \\
\midrule
\texttt{none} (default) & Uses sandbox workspace (\texttt{\textasciitilde{}/.openclaw/sandboxes}) \\
\texttt{ro} & Read-only mount of agent workspace \\
\texttt{rw} & Read-write mount of agent workspace \\
\bottomrule
\end{tabular}
\end{center}

\section{Minimal Enable Example}

\begin{lstlisting}[language=json]
{
  agents: {
    defaults: {
      sandbox: {
        mode: "all",
        scope: "session",
        workspaceAccess: "none",
        docker: {
          image: "openclaw-sandbox:bookworm-slim"
        }
      }
    }
  }
}
\end{lstlisting}

\section{Image Build}

\begin{lstlisting}[language=bash]
# Build sandbox image
scripts/sandbox-setup.sh

# Build sandbox browser image (optional)
scripts/sandbox-browser-setup.sh
\end{lstlisting}

The default sandbox container has \textbf{no network}; override with \texttt{docker.network}.

\section{Custom Bind Mounts}

Mount additional host directories into the container:

\begin{lstlisting}[language=json]
{
  agents: {
    defaults: {
      sandbox: {
        docker: {
          binds: [
            "/home/user/source:/source:ro",
            "/var/run/docker.sock:/var/run/docker.sock"
          ]
        }
      }
    }
  }
}
\end{lstlisting}

\section{setupCommand}

\texttt{setupCommand} runs \textbf{once} after container creation (not per run), used to install additional dependencies:

\begin{lstlisting}[language=json]
{
  agents: {
    defaults: {
      sandbox: {
        docker: {
          network: "bridge",
          setupCommand: "apt-get update && apt-get install -y python3"
        }
      }
    }
  }
}
\end{lstlisting}

\textbf{Note}: Package installation requires network egress (\texttt{network: "bridge"}) and a writable root filesystem.

\section{Multi-Agent Override}

Each agent can independently override sandbox configuration:

\begin{lstlisting}[language=json]
{
  agents: {
    list: [{
      id: "build",
      sandbox: {
        mode: "all",
        docker: {
          binds: ["/mnt/cache:/cache:rw"]
        }
      }
    }]
  }
}
\end{lstlisting}

\section{Debugging}

\begin{lstlisting}[language=bash]
openclaw sandbox explain  # Check effective sandbox config
\end{lstlisting}
